% from ega4-2-fr.tex
\oldpage[IV-2]{18}
donc $Z=f(F)$ est pro-constructible dans $Y$ (1.9.5, (vii)), et le corollaire
résulte de ce que $F$ est fermé.

On peut encore écrire le résultat de (2.3.11) sous la forme
$F=f^{-1}(f(X)\cap\overline{f(F)})$, autrement dit, les ensembles fermés du
sous-espace $f(X)$ de $Y$ sont les parties $Z\subset f(X)$ telles que
$f^{-1}(Z)$ soit fermé dans $X$~; cela signifie aussi que \emph{la topologie
induite par $Y$ sur $f(X)$ est quotient de celle de $X$ par la relation
d'équivalence définie par $f$}. En particulier~:

\begin{corollary}[2.3.12]
\label{IV.2.3.12-fr}
Soit $f:X\to Y$ un morphisme quasi-fidèlement plat (2.3.3) et quasi-compact.
Alors la topologie de $Y$ est quotient de celle de $X$ par la relation
d'équivalence définie par $f$ (autrement dit, pour que $Z\subset Y$ soit
ouvert (resp. fermé) dans $Y$, il faut et il suffit que $f^{-1}(Z)$ soit
ouvert (resp. fermé) dans $X$).
\end{corollary}

En effet, on a alors $f(X)=Y$.

\begin{corollary}[2.3.13]
\label{IV.2.3.13-fr}
Soient $X$, $Y$ deux $S$-préschémas, $f:X\to Y$ un $S$-morphisme fidèlement
plat et quasi-compact. Pour que $Y$ soit séparé au-dessus de $S$, il faut et il
suffit que l'immersion canonique
$j:X\times_Y X\to X\times_S X$ (I, 5.3.10) soit \emph{fermée}.
\end{corollary}

Notons pour cela que l'on a le diagramme commutatif (I, 5.3.5)
\[
\xymatrix{
  X\times_Y X \ar[r]^j \ar[d]_{\pi} &
  X\times_S X \ar[d]^{f\times_S f}\\
  Y \ar[r]_{\Delta_Y} & Y\times_S Y
}
\]
identifiant $X\times_Y X$ au produit des $(Y\times_S Y)$-préschémas $Y$ et
$X\times_S X$. Comme $f$ est surjectif, il en est de même de $\pi$ et de
$f\times_S f$, donc la diagonale $\Delta_Y(Y)$ a pour image réciproque par
$f\times_S f$ l'image $j(X\times_Y X)$ (I, 3.4.8). Comme $f\times_S f$ est
fidèlement plat et quasi-compact (1.1.2 et 2.2.13), il suffit d'appliquer
(2.3.12) à ce morphisme.

\begin{corollary}[2.3.14]
\label{IV.2.3.14-fr}
Soit $f:X\to Y$ un morphisme quasi-fidèlement plat et quasi-compact, et soit
$Z$ une partie de $Y$. Pour que $Z$ soit une partie localement fermée
pro-constructible de $Y$, il faut et il suffit que $f^{-1}(Z)$ soit une partie
localement fermée pro-constructible de $X$.
\end{corollary}

On sait déjà (1.9.12) que, pour que $Z$ soit une partie pro-constructible de
$Y$, il faut et il suffit que $f^{-1}(Z)$ soit une partie pro-constructible de
$X$. La condition est évidemment nécessaire. Pour prouver qu'elle est
suffisante, considérons le sous-préschéma fermé réduit $Y_1$ de $Y$ ayant
$\overline Z$ pour espace sous-jacent et, soit $X_1$ son image réciproque par
$f$, qui a pour espace sous-jacent
$f^{-1}(\overline Z)=\overline{f^{-1}(Z)}$ en vertu de (2.3.10). Comme
$f^{-1}(Z)$ est localement fermé dans $X$, il est ouvert dans
$\overline{f^{-1}(Z)}$, donc dans $X_1$~; or, le morphisme $f_1:X_1\to Y_1$
déduit de $f$ par restriction est quasi-fidèlement plat et quasi-compact
(1.1.2 et 2.3.3)~; on déduit donc de (2.3.12) que $Z$ est ouvert dans
$\overline Z$, ce qui montre que $Z$ est localement fermé dans $Y$.

\begin{remark}[2.3.15]
\label{IV.2.3.15-fr}
Il ne suffit pas, dans (2.3.12), de supposer seulement que $f$ est fidèlement
plat. Par exemple, prenons pour $Y$ une courbe algébrique réduite irréductible
